\editionPage[\pdfbookmark[0]{Fiches outils}{bm-outils-debut-parent}]{outils-debut}{Fiche outil}{Mesures et formules utiles}{À consulter au fil des chapitres.}
\begin{blocEdition}{outils-debut}

\editionEtape{Contours et aires}
\begin{tabularx}{\linewidth}{@{}lYY@{}}\toprule
Figure & Périmètre & Aire\\\midrule
Rectangle & $2(L+\ell)$ & $L\ell$\\[2mm]
Carré & $4c$ & $c^2$\\[2mm]
Triangle & somme des côtés & base $\times$ hauteur associée / 2\\[2mm]
Parallélogramme & somme des côtés & base $\times$ hauteur associée\\[2mm]
Disque / cercle & $2\pi r$ pour le cercle & $\pi r^2$ pour le disque\\\bottomrule
\end{tabularx}
\editionEtape{Unités}
$1$ km = $1000$ m ; $1$ m = $100$ cm ; $1$ m$^2$ = $10000$ cm$^2$ ;
$1$ m$^3$ = $1000$ L ; $1$ L = $1000$ mL ; $1$ mL = $1$ cm$^3$.\par\medskip
Une durée $t$ en heures vaut $60t$ minutes. Une vitesse en km/h se
divise par 3,6 pour obtenir des m/s.
\editionEtape{Angles et triangles}
Angle plat : $180^\circ$ ; tour complet : $360^\circ$.
Dans un triangle, la somme des angles vaut $180^\circ$.
Un triangle isocèle a des angles à la base égaux ; un triangle
équilatéral a trois angles de $60^\circ$.
\editionRetenir{La hauteur d’un triangle est perpendiculaire à la base choisie. Un côté oblique n’est pas automatiquement une hauteur.}

\end{blocEdition}
\editionPage{outils-2}{Fiche outil}{Données : ce que montre chaque représentation}{À consulter au fil des chapitres.}
\begin{blocEdition}{outils-2}

\editionEtape{Exemple exact}
La série de pointures $38;38;39;40;40;40;42;42$ a pour effectifs
2, 1, 3, 2 aux valeurs 38, 39, 40, 42. Un tableau ou des bâtons
à ces abscisses permettent de lire les extrêmes : étendue $42-38=4$.
Un disque dont les secteurs portent ces mêmes valeurs et effectifs
permet aussi ce calcul ; les angles seraient $90^\circ$, $45^\circ$,
$135^\circ$, $90^\circ$.
\editionEtape{Exemple regroupé}
Pour des durées dans $[0;10[$, $[10;20[$, $[20;30[$ et $[30;40[$,
avec effectifs 4, 8, 6, 2, les hauteurs de l’histogramme sont
proportionnelles à ces effectifs. Pour le construire, placer les bornes
sur l’axe horizontal, choisir une graduation verticale et dessiner quatre
rectangles accolés de même largeur.\par\medskip
L’étendue exacte ne peut être retrouvée : le minimum est dans $[0;10[$
et le maximum dans $[30;40[$. On sait seulement que l’étendue est
strictement comprise entre 20 et 40 minutes. Les valeurs exactes ont été
perdues par le regroupement.
\editionEtape{Une estimation annoncée}
En remplaçant chaque classe par son centre (5, 15, 25, 35),
on estime la moyenne par $(4\times5+8\times15+6\times25+2\times35)/20=18$ min.
Cette valeur est une estimation, pas une moyenne exacte des observations inconnues.

\end{blocEdition}
\editionPage{outils-3}{Fiche outil}{Décomposer en facteurs premiers avec un algorithme}{À consulter au fil des chapitres.}
\begin{blocEdition}{outils-3}

\editionEtape{Une procédure en français structuré}
Choisir un entier $N\ge2$. Mettre reste à N et diviseur à 2.
Tant que reste est supérieur à 1 :\par
\begin{itemize}
\item si le reste de la division de reste par diviseur vaut 0, afficher
diviseur et remplacer reste par reste / diviseur ;
\item sinon, augmenter diviseur de 1.
\end{itemize}
La répétition teste donc de nouveau le même diviseur après une division réussie.
En blocs, employer « répéter jusqu’à reste = 1 » et une condition à deux branches.
\editionEtape{Exemple pour 84}
On affiche 2, puis 2, puis 3, puis 7. Les quotients successifs après
division sont 42, 21, 7 et 1. La décomposition vaut
$84=2^2\times3\times7$. Un diviseur composé ne sera pas affiché :
ses facteurs premiers auraient été retirés auparavant.
\editionEtape{Contrôles}
Multiplier les facteurs affichés pour retrouver N.
Tester un nombre premier, par exemple 13 : un seul facteur est affiché.
La procédure doit refuser 0, 1 et les entrées non entières dans ce cadre.
Elle termine, car soit le reste diminue par une division, soit le diviseur
augmente jusqu’à atteindre un diviseur du reste.

\end{blocEdition}
\editionPage{outils-4}{Fiche outil}{Rédiger un raisonnement géométrique}{À consulter au fil des chapitres.}
\begin{blocEdition}{outils-4}

\editionEtape{Une longueur avec Pythagore}
« Le triangle ABC est rectangle en A. D’après le théorème de Pythagore,
$BC^2=AB^2+AC^2$. Donc… Comme BC est une longueur positive, BC = … ».
Ne pas remplacer le carré par la longueur avant la racine.
\editionEtape{Un parallélisme avec Thalès}
« A, M, B sont alignés dans cet ordre, ainsi que A, N, C.
On compare $AM/AB$ et $AN/AC$ : … Les deux rapports sont égaux.
La réciproque de Thalès permet de conclure que $(MN)\parallel(BC)$. »
\editionEtape{Un angle avec trigonométrie}
« Dans ABC rectangle en A, pour l’angle B, le côté AC est opposé et
AB adjacent. Donc $\tan B=AC/AB$. La fonction inverse donne
$B\approx\ldots^\circ$. » Nommer l’angle évite un rapport inversé.
\editionEtape{Une transformation}
Préciser son type, son centre ou son axe, puis son angle, son déplacement
ou son rapport. Pour une homothétie négative, dire explicitement que
le point image se situe de l’autre côté du centre.
\editionRetenir{Une rédaction courte suffit si elle contient les hypothèses utiles, la propriété, le calcul et la conclusion.}

\end{blocEdition}
